\documentclass[../../../include/open-logic-section]{subfiles}

\begin{document}

\olfileid{sth}{card-arithmetic}{fix}
\olsection{\texorpdfstring{$\aleph$-不动点}{aleph-不动点}}

在 \olref[spine][]{chap} 中，我们曾提到替代公理需要得到辩护，因为它迫使层叠变得相当高大。在进行了一些基数算术的讨论后，我们可以稍作展示，看看这个层叠有多高。

显然 $0 < \aleph_0$，且 $1 < \aleph_1$，且 $2 < \aleph_2$\ldots
而且确实，这种大小差异随每一步只会变得\emph{更大}。因此很容易猜想，对每个序数 $\kappa$，都有 $\kappa< \aleph_\kappa$。

但在 $\ZFC$ 下，这个猜想是\emph{不成立的}。事实上，我们可以证明存在\emph{$\aleph$-不动点}，即满足 $\kappa=\aleph_\kappa$ 的基数 $\kappa$。

\begin{prop}\ollabel{alephfixed}
存在一个 $\aleph$-不动点。
\end{prop}

\begin{proof}
使用递归定义：
\begin{align*}
	\kappa_0 &= 0\\
	\kappa_{n+1} &= \aleph_{\kappa_n}\\
	\kappa&= \bigcup_{n < \omega}\kappa_n
\end{align*}
根据 \olref[cardinals][classing]{unioncardinalscardinal}，$\kappa$ 是一个基数。现在有：
\[
	\kappa= \bigcup_{n < \omega} \kappa_{n+1} =
	\bigcup_{n < \omega}\aleph_{\kappa_n} =
	\bigcup_{\alpha < \kappa}\aleph_\alpha = \aleph_\kappa
\]
%	By construction, $\kappa$ is the least cardinal greater than each
%	$\kappa_n$. So $\aleph_\kappa$ is the least cardinal greater than
%	each $\aleph_{\kappa_n}$, and hence greater than each
%	$\kappa_{n+1}$. But equally $\kappa$ is the least cardinal greater
%	than each $\kappa_{n+1} = \aleph_{\kappa_n}$. So $\kappa=
%	\aleph_\kappa$.
\end{proof}


Boolos 曾撰文专门讨论我们刚构造的这个 $\aleph$-不动点。在文章开头指出 $\kappa$ 的存在之后，他写道：

\begin{quote}
    [$\kappa$ 是] 一个\emph{相当大的}数，在那些从未接触过集合论的人看来；大得似乎足以让任何断言至少存在 $\kappa$ 个对象的理论的真理性受到质疑。
    \citep[p.~257]{Boolos2000}
\end{quote}

而他最终在论文结尾这样发问：

\begin{quote}
    [我们是否] 怀疑，无论故事的开端怎样，到了这一步时，\emph{车轮只是在空转}，我们听到的已不再是关于任何\emph{实际情况}的描述？
    \citep[p.~268]{Boolos2000}
\end{quote}

如果我们确实已经脱离了“任何实际情况”，那么我们必须将指责的矛头直接指向替代公理。因为正是这条公理使得我们的证明得以成立。在这种情况下，人们大概会认为，Boolos 需要重新审视他几十年前提出的那个主张，即替代公理“\emph{没有}不良后果”（参见 \olref[replacement][extrinsic]{sec}）。

但 $\kappa$ 的存在真的如此糟糕吗？在此，考虑一下 Russell 的\emph{《项狄传》悖论}或许有所帮助。特里斯特拉姆·项狄用日记记录他的生活，但记录一天要花掉他一年的时间。随着每一年过去，项狄落后得越来越远：一年后，他只记录了一天，却已经度过了 364 天未记录的日子；两年后，他只记录了两天，却度过了 728 天未记录的日子；三年后，他只记录了三天，却度过了 1092 天未记录的日子……\footnote{忽略闰年。} 然而，如果项狄\emph{永生不死}，他将能够记录下每一天，因为他会在自己生命的第 $n$ 年记录第 $n$ 天。这样，“在时间的尽头”，项狄将拥有一本完整的日记。

那么：为什么这与下面这种想法如此不同呢？——即 $\alpha$ 小于 $\aleph_\alpha$（而且确实，越来越小，甚至令人绝望地小），直至 $\kappa$，在此时我们终于“赶上了”，使得 $\kappa
= \aleph_\kappa$？

抛开这个不谈，假设我们接受 $\ZFC$，让我们以一些关于不动点构造的趣味探讨来结束本节。直观上说，以下三个结果表明，存在一个（非平凡的）点，在该点处层叠的宽度等于其高度：

\begin{prop}\ollabel{bethfixed}
存在一个 $\beth$-不动点，即一个满足 $\kappa=
\beth_\kappa$ 的 $\kappa$。
\end{prop}

\begin{proof}
如 \olref{alephfixed}，用“$\beth$”代替“$\aleph$”。
\end{proof}

\begin{prop}\ollabel{stagesize}
$\card{V_{\omega+\alpha}} = \beth_{\alpha}$。若 $\omega \ordtimes
\omega \leq \alpha$，则 $\card{V_\alpha} = \beth_\alpha$。
\end{prop}

\begin{proof}
第一个断言可由简单的超限归纳法得出。
第二个断言随之成立，因为如果 $\omega \ordtimes \omega \leq \alpha$，则有 $\omega + \alpha = \alpha$。为了证明这一点，我们使用\olref[ord-arithmetic][]{chap}中关于序数算术的若干事实。
首先注意到 $\omega \ordtimes \omega = \omega \ordtimes (1 \ordplus \omega) = (\omega  \ordtimes 1) \ordplus (\omega\ordtimes\omega) = \omega \ordplus (\omega \ordtimes \omega)$。
现在，如果 $\omega \ordtimes \omega \leq \alpha$，即对某个 $\beta$ 有 $\alpha = (\omega\ordtimes\omega) \ordplus \beta$，那么 $\omega \ordplus \alpha = \omega \ordplus ((\omega \ordtimes \omega) \ordplus \beta) = (\omega \ordplus (\omega \ordtimes \omega)) \ordplus \beta = (\omega \ordtimes \omega) \ordplus \beta = \alpha$。
\end{proof}

\begin{cor}
存在 $\kappa$ 使得 $\card{V_\kappa} = \kappa$。
\end{cor}

\begin{proof}
令 $\kappa$ 为如 \olref{bethfixed} 所给出的一个 $\beth$-不动点。
显然 $\omega \ordtimes \omega < \kappa$。因此由 \olref{stagesize}，$\card{V_\kappa} =
\beth_\kappa= \kappa$。
\end{proof}


在 $V_\kappa$ 之下的阶段数，和 $V_\kappa$ 的元素个数一样多。
直观上说，$V_\kappa$ 的宽度和它的高度相等。
这非常具有《项狄传》的风格：我们通过取\emph{幂集}从一个阶段推进到下一个阶段，
从而每一步都让我们的层叠变得\emph{更}大。
但是，“到最后”，即在阶段 $\kappa$，层叠的宽度赶上了它的高度。

人们可能会问：\emph{层叠宽度匹配其高度的现象多久发生一次？}
答案是：\emph{和序数一样多。}但这需要稍加解释。

我们如下定义一个项 $\tau$。对任意 $A$，令：
\begin{align*}
	\tau_0(A) & \defis \card{A}\\
	\tau_{n+1}(A) & \defis \beth_{\tau_n(A)}\\
	\tau(A) & \defis \bigcup_{n < \omega}\tau_n(A)
\intertext{与 \olref{bethfixed} 类似，对任意 $A$，$\tau(A)$ 都是一个 $\beth$-不动点，并且显然有 $\card{A} < \tau(A)$。
现在考虑如下的递归定义：}
	W_0 &\defis 0\\
	W_{\alpha + 1} & \defis \tau(W_\alpha)\\
	W_\alpha & \defis \bigcup_{\beta < \alpha} W_\beta
	\text{，当 $\alpha$ 是极限序数时}
\end{align*}
该构造对所有序数都有定义。直观上说，$W$ 是从序数到 $\beth$-不动点的“一个单射”。
并且，和之前完全一样，对任意 $\alpha$，$V_{W_\alpha}$ 的宽度与高度相等。

\end{document}
